\chapter{Planteamiento del problema}\label{chap:problema}

\section{Contexto y brecha diagnosticada}
El diagnóstico de anemia falciforme en entornos de bajos recursos enfrenta tensiones entre exactitud clínica y viabilidad operativa. Técnicas consolidadas como la cromatografía líquida de alta resolución, la electroforesis y los ensayos moleculares exhiben sensibilidad y especificidad elevadas, pero su despliegue exige infraestructura especializada, consumibles costosos y personal altamente entrenado, lo que limita la cobertura de programas de cribado en regiones con mayor carga de enfermedad \cite{McGann2017,Shrestha2025,Arishi2021}. En ese panorama, la microscopía holográfica digital sin lentes (DLHM) emerge como una alternativa atractiva por su hardware simple y de bajo costo, capaz de capturar información rica en fase y amplitud; sin embargo, la clasificación automática directamente desde hologramas—sin reconstrucción numérica—sigue siendo un reto abierto que demanda metodologías de procesamiento y aprendizaje robustas \cite{Amann2019,BuitragoDuque2023,OConnor2020}.

\section{Pregunta de investigación}
A partir de esta brecha se define la pregunta que orienta el proyecto: \emph{¿es posible desarrollar un sistema de clasificación basado en hologramas DLHM y aprendizaje automático que pueda diferenciar entre glóbulos rojos sanos y falciformes a partir de sus hologramas sin reconstruir?} Esta formulación delimita explícitamente el dominio (eritrocitos), la modalidad de adquisición (DLHM) y la restricción metodológica clave (ausencia de reconstrucción), y establece el objetivo de alcanzar desempeño clínicamente útil en un flujo de trabajo computacional eficiente \cite{Amann2019,BuitragoDuque2023,OConnor2020}.

\section{Objetivos}
El propósito general es \emph{desarrollar un sistema de clasificación de hologramas de microscopía sin lentes para glóbulos rojos, basado en técnicas de aprendizaje automático, orientado a la detección de anemia falciforme} \cite{Amann2019,BuitragoDuque2023,OConnor2020}. Para concretarlo, el trabajo se estructura en tres frentes integrados. Primero, se establece un \textit{pipeline} de preprocesamiento para hologramas DLHM que comprende conversión a escala de grises, redimensionamiento a 
512
×
512
512×512 píxeles, normalización de intensidad y realce de contraste mediante técnicas locales (p.,ej., tipo CLAHE), incorporando filtrado en frecuencia cuando sea pertinente con el fin de optimizar la calidad de entrada para los modelos \cite{Amann2019,BuitragoDuque2023}. Luego, se evalúa el desempeño de diversas arquitecturas de aprendizaje automático en el escenario de clasificación \emph{sin reconstrucción}, adoptando un enfoque clásico sustentado en la extracción de descriptores de textura, forma y contenido frecuencial del holograma, y comparando clasificadores representativos bajo un protocolo experimental controlado \cite{OConnor2020,BuitragoDuque2023}. Finalmente, se valida rigurosamente el sistema mediante esquemas de validación cruzada anidada y procedimientos estadísticos que reduzcan el sesgo por tamaño muestral y eviten \emph{overfitting}, reportando métricas estándar de evaluación (área bajo la curva ROC, precisión balanceada, sensibilidad y especificidad) con estimaciones de incertidumbre reproducibles \cite{Vabalas2019,Tsamardinos2022}.

\section{Alcance y delimitaciones}
El estudio se circunscribe a la tarea de clasificación binaria de hologramas DLHM de eritrocitos en dos categorías—sanos y falciformes—sin realizar reconstrucción digital del campo complejo. No se aborda el diseño ni la optimización del montaje óptico, la segmentación celular individual, la integración con historias clínicas electrónicas ni la evaluación de patologías hematológicas distintas a la anemia falciforme. El desarrollo computacional se implementa en Python y se orienta a ejecutarse en recursos de cómputo moderados, priorizando la reproducibilidad experimental y la trazabilidad del flujo de datos \cite{Amann2019,BuitragoDuque2023,OConnor2020}.

\section{Limitaciones y riesgos}
Entre las posibles fuentes de riesgo se consideran la disponibilidad limitada de datos anotados, la heterogeneidad biológica y técnica de las muestras y la sensibilidad del desempeño a la selección de hiperparámetros. Estas circunstancias pueden inducir estimaciones optimistas si no se controlan adecuadamente. Para mitigarlas, se adoptan protocolos de validación estadísticamente sólidos—en particular, validación cruzada anidada y análisis de estabilidad—y se documentan todas las decisiones de preprocesamiento y modelado mediante reportes ejecutables que faciliten auditoría y réplica independiente \cite{Vabalas2019,Tsamardinos2022}. La generalización fuera del dominio inmediato de adquisición se reconoce como un objetivo de trabajos posteriores que requerirá cohortes externas y estudios multicentro \cite{BuitragoDuque2023,OConnor2020}.

\section{Hipótesis y productos esperados}
La hipótesis de trabajo sostiene que un flujo de procesamiento basado en descriptores de textura, morfología y frecuencia, combinado con clasificadores clásicos entrenados y evaluados bajo validación cruzada anidada, puede discriminar con fiabilidad hologramas DLHM de eritrocitos sanos y falciformes sin necesidad de reconstrucción, alcanzando desempeño competitivo con la literatura y con incertidumbre cuantificada \cite{Amann2019,BuitragoDuque2023,OConnor2020,Vabalas2019,Tsamardinos2022}. Como productos se esperan un conjunto de \emph{scripts} y modelos serializados listos para reproducción, bitácoras y reportes automatizados de métricas, y una guía operativa para su evaluación técnica posterior, en concordancia con las directrices institucionales de propiedad intelectual \cite{EAFIT2018}.